% ─────────────────────────────────────────────────────────────────────────────
%  Capítulo 3 – Frontend: Análisis léxico-sintáctico y AST
% ─────────────────────────────────────────────────────────────────────────────
\chapter{Frontend: An\'alisis l\'exico-sint\'actico}
\label{chap:frontend}

\section{An\'alisis l\'exico}

El analizador l\'exico se implementa en \texttt{src/lexer.l} mediante Flex.
Convierte el flujo de caracteres de entrada en una secuencia de \emph{tokens}
que el parser consumir\'a.

\subsection{Categor\'ias de tokens}

\begin{table}[H]
  \centering
  \caption{Categor\'ias de tokens reconocidas por el lexer}
  \label{tab:tokens}
  \begin{tabular}{llp{6cm}}
    \toprule
    \textbf{Categor\'ia} & \textbf{Patr\'on} & \textbf{Ejemplos} \\
    \midrule
    Palabras reservadas & exactas       & \texttt{int}, \texttt{while}, \texttt{return} \\
    Tipos               & exactas       & \texttt{int}, \texttt{double}, \texttt{boolean} \\
    Anotaciones formales& exactas       & \texttt{@pre}, \texttt{@post}, \texttt{@invariant} \\
    Identificadores     & \texttt{[a-zA-Z][a-zA-Z0-9]*} & \texttt{x}, \texttt{sum}, \texttt{absValue} \\
    Literales enteros   & \texttt{[0-9]+}               & \texttt{0}, \texttt{42} \\
    Literales reales    & \texttt{[0-9]+\textbackslash.[0-9]+} & \texttt{3.14} \\
    Operadores          & s\'imbolos    & \texttt{+}, \texttt{-}, \texttt{==}, \texttt{==>} \\
    Puntuaci\'on        & s\'imbolos    & \texttt{\{}, \texttt{;}, \texttt{(} \\
    \bottomrule
  \end{tabular}
\end{table}

\subsection{Tokens especiales para verificaci\'on}

Para la fase 5 se a\~nadieron tokens dedicados a las anotaciones de
especificaci\'on formal, colocados \emph{antes} de la regla general de
identificadores para que tengan precedencia:

\begin{lstlisting}[style=cpp, caption={Reglas l\'exicas para anotaciones formales}]
"result"      { RET_ATTR(result_token); }
"@pre"        { RET_ATTR(pre_token); }
"@post"       { RET_ATTR(post_token); }
"@invariant"  { RET_ATTR(invariant_token); }
"@variant"    { RET_ATTR(variant_token); }
"==>"         { RET_ATTR(implies_token); }
\end{lstlisting}

\section{An\'alisis sint\'actico}

El parser se implementa en \texttt{src/parser.y} con GNU Bison (gram\'atica
LALR(1)). Construye el AST de forma dirigida por la sintaxis mediante
acciones sem\'anticas en C++.

\subsection{Estructura gramatical principal}

\begin{lstlisting}[style=pseudo, caption={Reglas principales de la gram\'atica Bison}]
// Un programa es una clase con miembros
Program : public class id '(' ')' '{' Members '}'
Members : Members Member | Members OptPre | Members OptPost | /* empty */
Member  : Main | Method | DVar

// Metodo estatico sin parametros
Method  : public static Tipo id '(' ')' Bloque
        | public static void id '(' ')' Bloque

// Instrucciones
Instr   : DVar | Asig | Print | Return
        | if '(' Expr ')' Instr
        | if '(' Expr ')' Instr else Instr
        | while '(' Expr ')' OptInvariant OptVariant Instr
        | Bloque

// Anotaciones de especificacion formal
OptPre      : '@pre'       PredExpr | /* empty */
OptPost     : '@post'      PredExpr | /* empty */
OptInvariant: '@invariant' PredExpr | /* empty */
OptVariant  : '@variant'   PredExpr | /* empty */
PredExpr    : Expr '==>' Expr | Expr
\end{lstlisting}

\subsection{Gesti\'on de anotaciones formales}

Las anotaciones \texttt{@pre}/\texttt{@post} se acumulan en variables globales
\texttt{g\_pre} y \texttt{g\_post} antes de que se procese el m\'etodo. Cuando
se reduce una regla \texttt{Method}, el nodo \texttt{MethodNode} recibe las
precondici\'on y postcondici\'on almacenadas:

\begin{lstlisting}[style=cpp, caption={Vinculaci\'on de anotaciones al MethodNode}]
Method : public_token static_token Tipo id pari pard Bloque {
    MethodNode* m = new MethodNode($4->lexema, $3->tipo, $7, ...);
    m->precondition.reset(g_pre);   // @pre almacenada globalmente
    m->postcondition.reset(g_post); // @post almacenada globalmente
    g_pre = g_post = nullptr;
    $$ = m;
}
\end{lstlisting}

\section{\'Arbol Sint\'actico Abstracto (AST)}

El AST se representa mediante una jerarqu\'ia de clases C++ en
\texttt{include/ASTNodes.h}.

\subsection{Jerarqu\'ia de nodos}

\begin{figure}[H]
\centering
\begin{tikzpicture}[
  cls/.style={rectangle, draw, rounded corners=2pt,
              minimum width=3cm, minimum height=0.6cm,
              align=center, font=\small\ttfamily, fill=blue!8},
  arr/.style={-Stealth, thick},
  node distance=0.5cm and 1.2cm
]
\node[cls] (node)   {Node};
\node[cls, below left=0.8cm and 1.5cm of node]  (stmt)  {StmtNode};
\node[cls, below right=0.8cm and 1.5cm of node] (expr)  {ExprNode};

\node[cls, below=0.5cm of stmt] (assign) {AssignNode};
\node[cls, below=0.5cm of assign] (block)  {BlockNode};
\node[cls, below=0.5cm of block]  (ifn)    {IfNode};
\node[cls, below=0.5cm of ifn]    (whi)    {WhileNode};
\node[cls, below=0.5cm of whi]    (met)    {MethodNode};
\node[cls, below=0.5cm of met]    (ret)    {ReturnNode};

\node[cls, below=0.5cm of expr]   (bin)    {BinaryExprNode};
\node[cls, below=0.5cm of bin]    (una)    {UnaryExprNode};
\node[cls, below=0.5cm of una]    (ide)    {IdentifierNode};
\node[cls, below=0.5cm of ide]    (ilit)   {IntLiteralNode};
\node[cls, below=0.5cm of ilit]   (flo)    {FloatLiteralNode};
\node[cls, below=0.5cm of flo]    (cal)    {CallNode};
\node[cls, below=0.5cm of cal]    (res)    {ResultNode};

\draw[arr] (stmt) -- (node);
\draw[arr] (expr) -- (node);
\draw[arr] (assign) -- (stmt);
\draw[arr] (block) -- (stmt);
\draw[arr] (ifn) -- (stmt);
\draw[arr] (whi) -- (stmt);
\draw[arr] (met) -- (stmt);
\draw[arr] (ret) -- (stmt);
\draw[arr] (bin) -- (expr);
\draw[arr] (una) -- (expr);
\draw[arr] (ide) -- (expr);
\draw[arr] (ilit) -- (expr);
\draw[arr] (flo) -- (expr);
\draw[arr] (cal) -- (expr);
\draw[arr] (res) -- (expr);
\end{tikzpicture}
\caption{Jerarqu\'ia de nodos del AST (flecha = hereda de).}
\label{fig:ast}
\end{figure}

\subsection{Nodos especiales para verificaci\'on}

\begin{description}
  \item[\texttt{ResultNode}] Representa la palabra reservada \texttt{result}
    en las postcondiciones. Hace referencia al valor de retorno del m\'etodo sin
    necesidad de nombrarlo expl\'icitamente.
  \item[\texttt{MethodNode.precondition / .postcondition}] Campos
    \texttt{unique\_ptr<ExprNode>} que almacenan las expresiones de las
    anotaciones \texttt{@pre} y \texttt{@post}.
  \item[\texttt{WhileNode.invariant / .variant}] Campos an\'alogos para las
    anotaciones \texttt{@invariant} y \texttt{@variant} del bucle.
\end{description}

\begin{lstlisting}[style=cpp, caption={Nodo WhileNode con soporte de invariante}]
class WhileNode : public StmtNode {
public:
    std::unique_ptr<ExprNode> condition;
    std::unique_ptr<StmtNode> body;
    std::unique_ptr<ExprNode> invariant;  // @invariant annotation
    std::unique_ptr<ExprNode> variant;    // @variant annotation
    ...
};
\end{lstlisting}
